\documentclass[11pt]{article}
\usepackage[a4paper, margin=1in]{geometry}
\usepackage{hyperref}
\usepackage{titlesec}
\usepackage{enumitem}
\usepackage{listings}
\usepackage{xcolor}
\usepackage{tcolorbox}
\usepackage{amsmath}

\definecolor{codegray}{gray}{0.95}
\definecolor{darkblue}{rgb}{0.0,0.0,0.5}

\lstset{
  backgroundcolor=\color{codegray},
  basicstyle=\ttfamily\small,
  breaklines=true,
  frame=single,
  keywordstyle=\color{darkblue}\bfseries,
  commentstyle=\itshape\color{gray},
  stringstyle=\color{black}
}

\newtcolorbox{noteBox}{
  colback=blue!5,
  colframe=blue!60,
  title=\textbf{💡 Note},
  sharp corners,
  fonttitle=\bfseries
}

\newtcolorbox{warnBox}{
  colback=red!5,
  colframe=red!60,
  title=\textbf{⚠️ Warning},
  sharp corners,
  fonttitle=\bfseries
}

\newtcolorbox{exampleBox}{
  colback=green!5,
  colframe=green!60,
  title=\textbf{📝 Example},
  sharp corners,
  fonttitle=\bfseries
}

\title{\textbf{Complete Sui Move Revision Course}\\
\large From Zero to Confident: Comprehensive Guide\\
\large Abilities, Objects, OTW, Publisher, Capabilities, Option, Vectors, Loops, Events}
\author{Bootcamp Notes - Concise Version}
\date{}

\begin{document}

\maketitle
\tableofcontents
\newpage

% ============================================================
\section{Introduction: How to Use This Document}

\begin{noteBox}
This document covers all key concepts with concise explanations. Read explanations first, then study code examples.
\end{noteBox}

% ============================================================
\section{Basics: Move and Sui}

\subsection{What is a Smart Contract?}

A smart contract is a deterministic, verifiable, and secure program on a blockchain.

\begin{exampleBox}
\textbf{Analogy}: Like a vending machine that always gives the same product for the same input, and anyone can verify it worked correctly.
\end{exampleBox}

\subsection{What is Move?}

Move is a language designed for \textbf{assets}. Its philosophy: assets are special values that must be handled safely.

\begin{noteBox}
\textbf{Key Difference}: In Move, transferring an asset \textbf{moves} it (doesn't copy). This prevents accidental duplication.
\end{noteBox}

\subsection{What is Sui Move?}

Sui Move = Move + Sui's object system (ownership, unique IDs, fast transfers).

\begin{noteBox}
In Sui, almost everything important is an \textbf{object} with a unique ID and owner.
\end{noteBox}

\subsection{Modules, Packages, and Addresses}

\begin{itemize}
  \item \textbf{Package}: Collection of modules (what you publish)
  \item \textbf{Module}: Code file with functions (\texttt{package::module})
  \item \textbf{Address}: Unique identifier on-chain (Package ID)
\end{itemize}

% ============================================================
\section{Types and Structs}

\subsection{Basic Types}

\begin{itemize}
  \item \texttt{u64}: Unsigned 64-bit integer (0 to very large)
  \item \texttt{bool}: true/false
  \item \texttt{address}: Blockchain address
\end{itemize}

\subsection{Structs}

Structs are custom types you define:

\begin{lstlisting}[language=Rust]
module demo::basic {
    struct Point {
        x: u64,
        y: u64
    }
}
\end{lstlisting}

\begin{noteBox}
Structs are the main way to represent data in Move. In Sui, many structs become objects.
\end{noteBox}

% ============================================================
\section{Abilities (Extremely Important)}

\subsection{What are Abilities?}

Abilities are \textbf{permissions on a type} that control what you can do with values.

\begin{noteBox}
\textbf{Analogy}: Like file permissions - a type can be "read-only" or "modifiable" based on its abilities.
\end{noteBox}

\subsection{The 4 Abilities}

\subsubsection{1. \texttt{copy} - Permission to Duplicate}

\begin{warnBox}
\textbf{⚠️ Critical Rule}: Assets (money, NFTs, tokens) must \textbf{NEVER} have \texttt{copy}. Otherwise, you can duplicate money!
\end{warnBox}

\subsubsection{2. \texttt{drop} - Permission to Destroy}

Allows destroying or ignoring a value. If a value doesn't have \texttt{drop}, you must use it somewhere.

\subsubsection{3. \texttt{store} - Permission to Store}

Allows storing the value inside other structs. If a struct doesn't have \texttt{store}, you can't put it in another struct.

\subsubsection{4. \texttt{key} - Permission to be a Sui Object}

Allows a struct to exist as a top-level object on Sui with unique ID and owner.

\begin{noteBox}
\textbf{Rule}: If a struct has \texttt{key}, it \textbf{MUST} contain \texttt{id: UID}.
\end{noteBox}

\subsection{Why Abilities Matter}

\begin{lstlisting}[language=Rust]
// Safe data can have copy
struct Point has copy, drop, store {
    x: u64,
    y: u64
}

// Assets should NOT have copy
struct Coin has store {
    value: u64
}
// ❌ Cannot copy - prevents duplication
\end{lstlisting}

\subsection{Field Ability Rule}

A struct can only have an ability if \textbf{all its fields} support it.

% ============================================================
\section{Sui Objects and UID}

\subsection{What is an Object in Sui?}

A Sui object has:
\begin{enumerate}
  \item Unique ID (like a serial number)
  \item Owner (or is shared)
  \item Data fields
\end{enumerate}

\begin{noteBox}
\textbf{Analogy}: Like a physical book with ISBN (ID), owner, and pages (data).
\end{noteBox}

\subsection{What is UID?}

\texttt{UID} = Unique Identifier. Special type provided by Sui to guarantee unique object IDs.

\begin{lstlisting}[language=Rust]
struct MyObject has key {
    id: UID,      // Required for key
    value: u64
}
\end{lstlisting}

\subsection{Owned vs Shared Objects}

\begin{itemize}
  \item \textbf{Owned}: Belongs to one address
  \begin{itemize}
    \item Usually \textbf{fast} transactions (no global ordering needed)
  \end{itemize}
  
  \item \textbf{Shared}: Accessible by everyone
  \begin{itemize}
    \item \textbf{Slower} transactions (requires consensus ordering)
  \end{itemize}
\end{itemize}

\begin{warnBox}
If you put a huge vector in a shared object with many users mutating it, performance degrades (many conflicts).
\end{warnBox}

% ============================================================
\section{Publisher (Package Authorship Proof)}

\subsection{What is Publisher?}

\textbf{Publisher} is a special object created automatically when you publish a package. It proves: "I am the original author of this package."

\begin{noteBox}
\textbf{Analogy}: Like a certificate of authorship for your code.
\end{noteBox}

\subsection{Why it Exists}

Ensures only original creators can perform admin actions:
\begin{itemize}
  \item Restricted minting
  \item Updating rules
  \item Package upgrades (depending on settings)
\end{itemize}

\subsection{How it's Used}

\begin{lstlisting}[language=Rust]
public fun mint_nft(
    publisher: &Publisher,  // Must pass Publisher
    ctx: &mut TxContext
) {
    // Only creator (who has Publisher) can call this
}
\end{lstlisting}

\begin{noteBox}
If function requires \texttt{\&Publisher} and caller doesn't have it, transaction fails automatically.
\end{noteBox}

% ============================================================
\section{One-Time Witness (OTW)}

\subsection{The Problem OTW Solves}

Sometimes you must ensure initialization happens \textbf{exactly once}:
\begin{itemize}
  \item Creating a token type
  \item Creating global config
  \item Creating a unique registry
\end{itemize}

\subsection{What is an OTW?}

A One-Time Witness is a struct that:
\begin{enumerate}
  \item Is created once (at publish time)
  \item Is then consumed (used up)
  \item Cannot be recreated
\end{enumerate}

\begin{noteBox}
\textbf{Analogy}: Like a one-time concert ticket - once used, it's destroyed and can never be recreated.
\end{noteBox}

\subsection{Simple OTW Pattern}

\begin{lstlisting}[language=Rust]
module demo::otw {
    struct INIT {}  // Empty, no abilities = perfect OTW
    
    public fun init(witness: INIT, ctx: &mut TxContext) {
        // This can only run once because INIT is one-time
        // Setup your singleton/config/registry here
    }
}
\end{lstlisting}

\begin{warnBox}
If you give abilities (like \texttt{store} or \texttt{copy}) to the OTW struct, you break the "one-time" guarantee.
\end{warnBox}

% ============================================================
\section{Capabilities (Permission as Object)}

\subsection{What is a Capability?}

A \textbf{capability} is an \textbf{object} that represents permission.

\begin{noteBox}
\textbf{Paradigm Shift}: Instead of "if sender is admin", Sui style: "if you own AdminCap, you are admin."
\end{noteBox}

\subsection{Why Capabilities are Powerful}

\begin{itemize}
  \item \textbf{Unforgeable}: Only your module can create them
  \item \textbf{Transferable}: Can give admin to someone else
  \item \textbf{Storable in multisigs/DAOs}: Enables shared governance
\end{itemize}

\subsection{Beginner Example: MintCap}

\begin{lstlisting}[language=Rust]
module demo::capabilities {
    struct Coin has key {
        id: UID,
        value: u64
    }
    
    // Whoever owns this has mint authority
    struct MintCap has key {
        id: UID
    }
    
    public fun mint(
        cap: &MintCap,  // Must pass MintCap
        amount: u64,
        ctx: &mut TxContext
    ): Coin {
        Coin { id: object::new(ctx), value: amount }
    }
}
\end{lstlisting}

\begin{noteBox}
User must pass \texttt{\&MintCap}. If they don't own it, they can't supply it to the transaction.
\end{noteBox}

\subsection{Capabilities vs Publisher vs OTW}

\begin{itemize}
  \item \textbf{Publisher}: Proves package publisher (authorship/admin at package level)
  \item \textbf{OTW}: Proves initialization can only happen once
  \item \textbf{Capability}: Grants ongoing permission (mint, burn, admin update, etc.)
\end{itemize}

% ============================================================
\section{Option<T> (Maybe a Value, Maybe Nothing)}

\subsection{What is Option?}

Move avoids \texttt{null}. Instead uses:
\[
Option<T> = \text{Some}(T) \ \text{or} \ \text{None}
\]

\begin{noteBox}
\textbf{Why it matters}: Forces explicit handling - "Is there a value or not?"
\end{noteBox}

\subsection{Simple Example}

\begin{lstlisting}[language=Rust]
use std::option;

let a = option::some(10);      // Some(10)
let b = option::none<u64>();   // None

if (option::is_some(&a)) {
    let x = option::extract(a);  // Consumes Option
};

// To read without consuming
let y = *option::borrow(&b);  // Borrows, doesn't consume
\end{lstlisting}

\begin{noteBox}
\textbf{Key Difference}:
\begin{itemize}
  \item \texttt{extract}: \textbf{Consumes} Option (destroys it)
  \item \texttt{borrow}: \textbf{Borrows} Option (doesn't destroy it)
\end{itemize}
\end{noteBox}

\begin{warnBox}
If you \texttt{extract} on \texttt{None}, transaction fails. Always check with \texttt{is\_some()} first.
\end{warnBox}

% ============================================================
\section{Vectors (Dynamic Lists)}

\subsection{What is vector<T>?}

A vector is a dynamic list:
\[
vector<T> = [T, T, T, \dots]
\]

\subsection{Core Operations}

\begin{itemize}
  \item Create: \texttt{vector::empty<T>()}
  \item Add: \texttt{vector::push\_back(\&mut v, x)}
  \item Remove: \texttt{vector::pop\_back(\&mut v)}
  \item Length: \texttt{vector::length(\&v)}
  \item Read: \texttt{vector::borrow(\&v, i)}
\end{itemize}

\subsection{Beginner Example}

\begin{lstlisting}[language=Rust]
use std::vector;

let v = vector::empty<u64>();
vector::push_back(&mut v, 10);  // v = [10]
vector::push_back(&mut v, 20);  // v = [10, 20]

let len = vector::length(&v);           // 2
let first = *vector::borrow(&v, 0);     // 10
let last = vector::pop_back(&mut v);    // 20
\end{lstlisting}

\begin{noteBox}
\textbf{Important}:
\begin{itemize}
  \item \texttt{borrow} returns a reference - use \texttt{*} to get value
  \item \texttt{pop\_back} removes and returns last element
  \item Indices start at 0
\end{itemize}
\end{noteBox}

\begin{warnBox}
If you \texttt{borrow} with an index that doesn't exist, transaction aborts. Always check length first.
\end{warnBox}

\subsection{Sui Performance Note}

If a vector is in a shared object with many users mutating it, there will be conflicts (slower). Common solution: use multiple objects ("sharding") instead of one giant vector.

% ============================================================
\section{Loops (while, break, continue)}

\subsection{Move Loops are Explicit}

Move uses \texttt{while}. No classic \texttt{for} loop.

\begin{noteBox}
\textbf{Why while?} Forces explicit condition handling and makes code clearer.
\end{noteBox}

\subsection{Iterating a Vector}

\begin{lstlisting}[language=Rust]
let i = 0;
let len = vector::length(&v);

while (i < len) {
    let val = *vector::borrow(&v, i);
    // Do something with val
    i = i + 1;  // Move to next element
}
\end{lstlisting}

\subsection{break and continue}

\begin{lstlisting}[language=Rust]
while (i < len) {
    let val = *vector::borrow(&v, i);
    
    if (val == 0) { break; };      // Exit loop
    if (val == 1) { 
        i = i + 1;
        continue;                   // Skip to next iteration
    };
    
    i = i + 1;
}
\end{lstlisting}

\begin{noteBox}
\textbf{Explanation}:
\begin{itemize}
  \item \texttt{break}: Exit loop immediately
  \item \texttt{continue}: Skip to next iteration
  \item Don't forget to increment \texttt{i} before \texttt{continue}!
\end{itemize}
\end{noteBox}

\begin{warnBox}
Loops cost gas. Very large loops can make transactions expensive or fail due to gas limits.
\end{warnBox}

% ============================================================
\section{Events (Contract Communication)}

\subsection{What is an Event?}

An event is a \textbf{structured log} emitted during a transaction.

\begin{noteBox}
\textbf{Analogy}: Like a journal entry - records "what happened" but doesn't change persistent state.
\end{noteBox}

\subsection{Events vs Objects}

\begin{itemize}
  \item \textbf{Object}: Persistent on-chain state (can be read later)
  \item \textbf{Event}: Log output (useful off-chain, not used as state)
\end{itemize}

\subsection{Defining an Event}

\begin{lstlisting}[language=Rust]
struct MintEvent has copy, drop {
    owner: address,
    amount: u64
}
\end{lstlisting}

\begin{noteBox}
Events usually have \texttt{copy, drop} because they're not assets - copying is safe.
\end{noteBox}

\subsection{Emitting an Event}

\begin{lstlisting}[language=Rust]
use sui::event;

event::emit(MintEvent {
    owner: tx_context::sender(ctx),
    amount: 100
});
\end{lstlisting}

\begin{noteBox}
\textbf{How it works}:
\begin{itemize}
  \item \texttt{event::emit()} records event in transaction
  \item Frontends can listen to these events
  \item Excellent for real-time UI updates
\end{itemize}
\end{noteBox}

\subsection{When to Use Each}

\begin{itemize}
  \item \textbf{Use Object} if: Need to read/modify state later, information must persist
  \item \textbf{Use Event} if: Just want to notify something happened, frontend reactions, history log
\end{itemize}

% ============================================================
\section{Putting It All Together (Complete Pattern)}

\subsection{Goal}

Create a system with:
\begin{enumerate}
  \item Shared \textbf{Config} object created once (OTW)
  \item \textbf{AdminCap} capability to update it
  \item Events when updates happen
  \item Optional admin address (Option)
\end{enumerate}

\subsection{Complete Example}

\begin{lstlisting}[language=Rust]
module demo::app {
    use sui::object::{UID};
    use sui::tx_context::TxContext;
    use sui::transfer;
    use sui::event;
    use std::option::{Self, Option};

    struct INIT {}  // OTW

    struct AdminCap has key { id: UID }

    struct Config has key {
        id: UID,
        admin: Option<address>,  // Optional admin
        threshold: u64
    }

    struct ConfigUpdated has copy, drop {
        new_threshold: u64
    }

    // Initialize once
    public fun init(w: INIT, ctx: &mut TxContext) {
        let cap = AdminCap { id: object::new(ctx) };
        let cfg = Config {
            id: object::new(ctx),
            admin: option::some(tx_context::sender(ctx)),
            threshold: 0
        };

        transfer::transfer(cap, tx_context::sender(ctx));
        transfer::share_object(cfg);
    }

    // Only admin can update
    public fun set_threshold(
        cap: &AdminCap,
        cfg: &mut Config,
        v: u64
    ) {
        cfg.threshold = v;
        event::emit(ConfigUpdated { new_threshold: v });
    }
}
\end{lstlisting}

\begin{noteBox}
\textbf{Explanation}:
\begin{enumerate}
  \item \textbf{OTW (INIT)}: Guarantees \texttt{init()} runs only once
  \item \textbf{AdminCap}: Capability object for admin permission
  \item \textbf{Config}: Shared object with global configuration
  \item \textbf{Option<address>}: Optional admin (Some or None)
  \item \textbf{ConfigUpdated}: Event emitted on config change
  \item \texttt{init()}: Creates everything once, gives AdminCap to creator, shares Config
  \item \texttt{set\_threshold()}: Requires AdminCap, updates Config, emits event
\end{enumerate}
\end{noteBox}

% ============================================================
\section{Common Beginner Mistakes}

\begin{enumerate}
  \item \textbf{Giving assets \texttt{copy} ability} (DANGEROUS!)
  \begin{itemize}
    \item Solution: Never give \texttt{copy} to assets
  \end{itemize}
  
  \item \textbf{Forgetting \texttt{id: UID} in \texttt{has key} structs}
  \begin{itemize}
    \item Solution: Always include \texttt{id: UID} in Sui objects
  \end{itemize}
  
  \item \textbf{Treating events as persistent storage}
  \begin{itemize}
    \item Solution: Use objects for persistent state
  \end{itemize}
  
  \item \textbf{Writing huge loops on shared vectors}
  \begin{itemize}
    \item Solution: Use multiple objects (sharding) or limit size
  \end{itemize}
  
  \item \textbf{Not using capabilities for admin actions}
  \begin{itemize}
    \item Solution: Use capabilities for permissions
  \end{itemize}
  
  \item \textbf{Confusing OTW (one-time init) with Publisher (authorship proof)}
  \begin{itemize}
    \item Solution: Understand difference and use each correctly
  \end{itemize}
\end{enumerate}

% ============================================================
\section{Final Mental Model}

\begin{center}
\Large\textbf{Sui Move = Objects + Ownership + Abilities + Capabilities}
\end{center}

\subsection{Key Concepts}

\begin{enumerate}
  \item \textbf{Abilities}: Prevent unsafe behavior (copy money, drop assets)
  \item \textbf{Objects}: Persistent state (with UID)
  \item \textbf{OTW}: Ensures one-time initialization
  \item \textbf{Publisher}: Proves package authorship and protects admin actions
  \item \textbf{Capabilities}: "Permission objects"
  \item \textbf{Option}: Safe "maybe" values (no null)
  \item \textbf{Vector}: Stores lists, loops iterate them
  \item \textbf{Events}: Broadcast what happened to off-chain world
\end{enumerate}

\subsection{Sui Move Philosophy}

\begin{noteBox}
\textbf{Fundamental Principles}:
\begin{itemize}
  \item \textbf{Security by Design}: Language forces you to be explicit about assets
  \item \textbf{Objects as Permissions}: Owning an object = having permission
  \item \textbf{Explicit over Implicit}: Must say exactly what you want to do
  \item \textbf{Composition}: Build complex systems from simple building blocks
\end{itemize}
\end{noteBox}

\end{document}
